\documentclass{article}

% Pass 'final' to the neurips_2023 package to remove line numbers if needed
\usepackage[final]{neurips_2023}

% \usepackage[style=authoryear, backend=biber]{biblatex} 
% \addbibresource{main.bib}

\usepackage{graphicx}

\usepackage[utf8]{inputenc} % allow utf-8 input
\usepackage[T1]{fontenc}    % use 8-bit T1 fonts
\usepackage{hyperref}       % hyperlinks
\usepackage{url}            % simple URL typesetting
\usepackage{booktabs}       % professional-quality tables
\usepackage{amsfonts}       % blackboard math symbols
\usepackage{nicefrac}       % compact symbols for 1/2, etc.
\usepackage{microtype}      % microtypography
\usepackage{xcolor}         % colors

\usepackage{tabularx}
\usepackage{booktabs}
\usepackage{makecell}

\usepackage{hyperref}
\usepackage{cleveref}

\title{Summarization: BoxMap: Efficient Structural Mapping and Navigation}

% Paper BibTex
% @inproceedings{dennler2025contrastive,
%   title={Contrastive Learning from Exploratory Actions: Leveraging Natural Interactions for Preference Elicitation},
%   author={Dennler, Nathaniel and Nikolaidis, Stefanos and Matari{\'c}, Maja},
%   booktitle={2025 20th ACM/IEEE International Conference on Human-Robot Interaction (HRI)},
%   pages={778--788},
%   year={2025},
%   organization={IEEE}
% }

\author{
  Summarized by Xiaonan (Nice) Wang \thanks{wangxiaonannice@gmail.com} \\
  % Paper fr last-names (year)
  Summarization\thanks{\textbf{Disclaimer:} This summarization is for research and study purposes. It represents a personal interpretation and may contain inaccuracies. Feedback or corrections via email are highly appreciated.} Generated based on Paper from Wang et al. (2025) \\
  Topic: Human-Robot Interaction, Multimodal Trajectory Feature Generation, Preference Learning \\
}

\begin{document}

\maketitle

\begin{abstract}
This document summarizes the core contributions and methodology of the paper "BoxMap: Efficient Structural Mapping and Navigation, \cite{wang2025boxmap}", focusing on its' main ideas and the core blocks.
\end{abstract}

\section{Overview: Core Questions and Answers}

\paragraph{(1) What is the problem?} 
\begin{itemize}
  \item Scenario (the Overall Task): \textbf{Customize} Robot Behaviors that \textbf{Align with Human Preference} -> a.k.a \textit{Preference Learning}
  \item \textit{Low-Dimensional} \textbf{Preference-Aligned (Semantic)} Trajectory \textbf{Features Generation} using Feature Learning Model (a.k.a Feature Generator)
\end{itemize}

\paragraph{(2) Why need to solve this problem?}
\begin{itemize}
  \item Generating semantic human preference-aligned trajectory features when customizing a robot, in order to reduce user interactions for eliciting preferences and improve user satisfaction;
  \item Avoiding data-labelling process like proxy tasks before customizing the robot.
\end{itemize}

\paragraph{(3) How is it different from prev.?}
\begin{itemize}
  \item Compared with \textit{self-supervised methods} like AEs or VAEs which result in \textit{physically representative features only}, CLEA will learn user preference-aligned features;
  \item Compared with \textit{hand-crafted} and \textit{prox tasks}:
  \begin{itemize}
    \item \textbf{Complete Feature Spaces} (compared with hand-crafted)
    \item \textbf{No Extra Labelling Process} (since exploratory search just additionally generates needed data when achieving primary goal)
  \end{itemize}
\end{itemize}

\paragraph{(4) Why is it better than prev.? (Advantages)} 
\begin{itemize}
  \item \textbf{Better Semantic Alignment}
  \item \textbf{Lower Human Effort}
  \item \textbf{Trade-off:} Obtained semantic representations without extra human efforts.
  \item Robot learning as an \textbf{exploratory search} (instead of \textit{information retrieval}) interaction is useful for novice robot users (\textbf{complete}, \textbf{simple}, \textbf{minimal} and \textbf{explainable})
\end{itemize}

\paragraph{(5) What is the approach itself?}
\begin{itemize}
  \item \textbf{CLEA:}
  \begin{itemize}
    \item \textbf{Contrastive} Learning from \textbf{Exploratory} Actions
    \item A \textit{\textbf{Generalizable} Trajectory Feature Learning Framework} that provides \textit{Low-Dimensional Feature Vector Generator}
  \end{itemize}
\end{itemize}

\begin{table}[h]
\centering
\caption{Functional Roles of BoxMap Components}
\label{tab:boxmap_components}
  \begin{tabularx}{\textwidth}{l l l X}
    \toprule
    Component & Input & Primary Function & Core Logic and Value \\ \midrule
    CNN Backbone & Occupancy Grid Map & \makecell[l]{Dimensionality Reduction and \\ Local Feature Extraction} & (1) \textbf{Downsamples} inputs into lower-resolution semantic feature maps. \newline (2) Reduces pixel count to prevent $O(N^2)$ complexity explosion in Attention. \newline (3) Provides \textbf{Inductive Bias} (Translation Invariance). \\ \addlinespace
    Transformer Encoder & Feature Maps & Global Context Encoding & Uses \textbf{Self-attention} to model long-range dependencies between pixels (e.g., relating distant walls). \\ \addlinespace
    Object Queries & Learnable Vectors & Entity Proposals & Acts as $M$ placeholders (occupants) to query potential objects in the scene. \\ \addlinespace
    Transformer Decoder & \makecell[l]{Queries + \\ Encoded Features (via \textbf{Cross-attention})} & Object Manifestation & Maps queries to specific $M$ \textbf{Box Embeddings}. \\ \addlinespace
    Decoder Self-Attention & Intermediate Queries & Parallel Coordination & Unmasked Attention allows queries to communicate and avoid redundant detections (Parallel de-duplication). \\ \addlinespace
    Hierarchical Loss & Predicted Boxes vs. GT & Detail Enhancement & Subtracts large objects (rooms) from TSDF to focus on small topological details (doors). \\ \addlinespace
    Prediction Heads & Box Embeddings & Geometric Mapping & Projects embeddings into box coordinates and class labels (specifying existence). \\ \bottomrule
  \end{tabularx}
\end{table}


\paragraph{(6) What are the applications of it?} 
\begin{itemize}
  \item Human-Robot Interaction Field
  \item A Generalizable Trajectory Feature Learning Framework that can be used in Assistive/Guidance Robots, etc.
\end{itemize}

\section{The Structure}

\paragraph{Summarized Block-Diagram} 
% The summarized block-diagram of \textit{experimental workflow} see \cref{fig:clea}.

% \begin{figure}[tb]
%   \centering
%   \includegraphics[width=0.95\textwidth]{fig/clea_experimental_workflow_diagram.png}
%   \caption{Summarized Block-Diagram of CLEA Experimental Workflow}
%   \label{fig:clea}
% \end{figure}

\paragraph{Original in Paper}

% The original block-diagram see \cref{fig:original}.

% \begin{figure}[tb]
%   \centering
%   \includegraphics[width=0.95\textwidth]{fig/orig.png}
%   \caption{Original Block-Diagram in Paper}
%   \label{fig:original}
% \end{figure}

\section{Open Questions}

\begin{enumerate}
    \item \textbf{Purpose of Symmetric Triplet Loss:} 
    Does the second half of Eq. 3 make sense? Will it increase the risk of overfitting while having limited effectiveness in promoting convergence?
    \item \textbf{Generalizabilityity and Scalability:}
    Is the trained CLEA feature generator only can be integrated into the robot customization process for online inference to obtain feature vectors? How about other scenarios like utilizing CLEA generator offline to generate vectors to build a feature vector database? If so, does the trained generator meet the generalizabilityity quirements?
\end{enumerate}

\bibliographystyle{abbrvnat}
\bibliography{main}
\end{document}
